\documentclass[11pt, aspectratio=169]{beamer}

% ============ PACKAGES ============
\usepackage[utf8]{inputenc}
\usepackage[french]{babel}
\usepackage{amsmath, amssymb, amsfonts}
\usepackage{siunitx}
\usepackage{booktabs}
\usepackage{tabularx}
\usepackage{graphicx}
\usepackage{tikz}
\usepackage{xcolor}
\usepackage{multicol}
\usepackage{geometry}
\usepackage{hyperref}
\usepackage{tcolorbox}

% ============ THEME & COLORS ============
\usetheme{Madrid}
\usecolortheme{default}

% Couleurs personnalisées (modèle respecté)
\definecolor{deepblue}{RGB}{0, 51, 102}
\definecolor{darkgreen}{RGB}{34, 139, 34}
\definecolor{lightgray}{RGB}{240, 240, 240}
\definecolor{accent1}{RGB}{220, 53, 69}
\definecolor{accent2}{RGB}{40, 167, 69}

\setbeamercolor{structure}{fg=deepblue}
\setbeamercolor{title}{bg=deepblue, fg=white}
\setbeamercolor{frametitle}{bg=deepblue, fg=white}
\setbeamercolor{alerted text}{fg=accent1}

% ============ SETTINGS ============
\setbeamertemplate{navigation symbols}{}
\setbeamerfont{title}{size=\Large, series=\bfseries}
\setbeamerfont{frametitle}{size=\large, series=\bfseries}

% ============ DOCUMENT INFO ============
\title{Détection précoce du Diabète de Type 1 au Cameroun}
\subtitle{par Apprentissage Automatique : utilisation de biomarqueurs génétiques émergents}
\author{SORELLE PERELLE NGUIAKAM \\ Étudiante en Master 2 Physique/Biophysique}
\institute{Université de Yaoundé I \\ Département de Physique}
\date{Année académique 2025--2026}

\begin{document}

% ============ SLIDE 1: TITLE PAGE ============
\frame{\titlepage}

% ============ SLIDE 2: OUTLINE ============
\begin{frame}{Plan de la Présentation}
  \begin{enumerate}
    \item Diabète de Type 1 idiopathique en Afrique subsaharienne
    \item Biomarqueurs et approches diagnostiques traditionnelles
    \item Apprentissage automatique pour la détection du DT1
    \item Méthodologie computationnelle du projet
  \end{enumerate}
\end{frame}

% ============================================================================
% SECTION 1: DIABÈTE DE TYPE 1 IDIOPATHIQUE EN AFRIQUE SUBSAHARIENNE
% ============================================================================
\section{Diabète de Type 1 idiopathique en Afrique subsaharienne}

% ============ SLIDE 3 ============
\begin{frame}{Définitions et Classification du DT1}
  \textbf{Définition du DT1 [IDF Atlas 2021]} :
  
  \vspace{0.6em}
  
  \begin{columns}[T]
    \column{0.50\textwidth}
    \textbf{\color{accent2} DT1 Auto-immun} :
    \begin{itemize}
      \item Destruction cellules $\beta$ pancréatiques
      \item Présence auto-anticorps (GADA, IA-2A)
      \item Traitement : insuline à vie
      \item Forme classique Occident
    \end{itemize}
    
    \column{0.50\textwidth}
    \textbf{\color{accent1} DT1 Idiopathique} :
    \begin{itemize}
      \item Absence auto-anticorps détectables
      \item Mécanisme non élucidé
      \item 30--40\% cas africains vs 5--10\% Occident
      \item Diagnostic plus complexe
    \end{itemize}
  \end{columns}
  
  \vspace{1em}
  
  \textbf{Distinction fondamentale} : DT1 (déficit absolu insuline) $\neq$ DT2 (résistance insuline + déficit relatif)
\end{frame}

% ============ SLIDE 4 ============
\begin{frame}{Classification : DT1 / DT2 / Formes Hybrides}
  \begin{table}[h]
    \tiny
    \centering
    \begin{tabularx}{\textwidth}{|l|X|X|X|}
      \toprule
      \textbf{Critère} & \textbf{Type 1} & \textbf{Type 2} & \textbf{KPD (Ketosis-Prone)} \\
      \midrule
      \textbf{Étiologie} & Auto-immune ou idiopathique & Résistance insuline & Idiopathique africain \\
      \textbf{Âge} & $<40$ ans typique & $>40$ ans typique & Variable \\
      \textbf{Auto-anticorps} & Présents (80--95\%) & Absents & Absents (forme africaine) \\
      \textbf{Acidocétose} & Fréquente au diagnostic & Rare & Très fréquente \\
      \textbf{Rémission} & Jamais & N/A & Possible temporairement \\
      \textbf{Prévalence Afrique} & 10\% diabètes & 85\% diabètes & 5--10\% diabètes \\
      \bottomrule
    \end{tabularx}
  \end{table}
  
  \vspace{0.8em}
  
  \textbf{\textcolor{accent1}{Ketosis-Prone Diabetes (KPD)}} : Forme hybride atypique, fréquente en Afrique subsaharienne, présentation aiguë avec possibilité de rémissions partielles
\end{frame}

% ============ SLIDE 5 ============
\begin{frame}{Formes Idiopathiques et Spécificités Africaines}
  \textbf{\textcolor{accent1}{Proportion élevée de formes idiopathiques}} :
  
  \vspace{0.8em}
  
  \begin{columns}[T]
    \column{0.50\textwidth}
    \textbf{\color{accent2} Présentation clinique} :
    \begin{itemize}
      \item Début aigu et brutal
      \item Acidocétose diabétique (ACD) fréquente
      \item Rémissions partielles possibles (KPD)
      \item Réponse initiale variable
      \item Évolution imprévisible
    \end{itemize}
    
    \column{0.50\textwidth}
    \textbf{\color{accent1} Contraintes diagnostiques} :
    \begin{itemize}
      \item Tests auto-anticorps négatifs
      \item Coût tests immunologiques : \$50--150 USD
      \item Infrastructures laboratorires limitées
      \item Personnel formé insuffisant
      \item Diagnostic souvent tardif
    \end{itemize}
  \end{columns}
  
  \vspace{1em}
  
  \textcolor{accent1}{\textbf{Contexte ressources limitées}} : Diagnostic repose principalement sur clinique + glycémie capillaire
\end{frame}

% ============ SLIDE 6 ============
\begin{frame}{Épidémiologie du DT1 en Afrique subsaharienne}
  \textbf{Données globales et régionales [IDF 2020--2025]} :
  
  \vspace{0.8em}
  
  \begin{columns}[T]
    \column{0.50\textwidth}
    \textbf{\color{accent2} Afrique subsaharienne} :
    \begin{itemize}
      \item $\approx$ 0,8--1,0 million personnes
      \item $\approx$ 10\% cas mondiaux DT1
      \item Tendance croissante 2010--2025
      \item Surveillance sanitaire fragmentée
      \item Épidémiologie sous-documentée
      \item Registres nationaux quasi-inexistants
    \end{itemize}
    
    \column{0.50\textwidth}
    \textbf{\color{accent1} Focus Cameroun [Katte et al. 2023]} :
    \begin{itemize}
      \item Prévalence $<1\%$ population générale
      \item Cas pédiatriques : augmentation majeure
      \item \textbf{Mortalité 1ère année : 25--30\%}
      \item Diagnostic souvent en urgence (ACD)
      \item \textbf{Absence registre national}
      \item Données hospitalières fragmentaires
    \end{itemize}
  \end{columns}
\end{frame}

% ============ SLIDE 7 ============
\begin{frame}{Enjeux Cliniques et de Santé Publique}
  \begin{tcolorbox}[colback=accent1!15, title=\textbf{\small Conséquences du diagnostic tardif}]
    \begin{itemize}
      \item \textbf{Acidocétose diabétique (ACD)} : 50--70\% diagnostics vs 15--30\% Europe
      \item \textbf{Décès précoce} : 25--30\% mortalité 1ère année vs $<5\%$ Occident
      \item \textbf{Complications chroniques} : rétinopathie, néphropathie accélérées
      \item \textbf{Coût économique} : Insuline = 10--30\% revenu mensuel familial
      \item \textbf{Ruptures approvisionnement} : Fréquentes et imprévisibles
    \end{itemize}
  \end{tcolorbox}
  
  \vspace{0.8em}
  
  \textbf{\color{accent2} Impact sur enfants et adolescents} :
  \begin{itemize}
    \item Scolarité interrompue
    \item Stigmatisation sociale
    \item Qualité de vie fortement dégradée
  \end{itemize}
  
  \vspace{0.6em}
  
  \textbf{\textcolor{accent1}{Besoin urgent}} : Outils de détection précoce adaptés au contexte africain
\end{frame}

% ============================================================================
% SECTION 2: BIOMARQUEURS ET APPROCHES DIAGNOSTIQUES TRADITIONNELLES
% ============================================================================
\section{Biomarqueurs et approches diagnostiques traditionnelles}

% ============ SLIDE 8 ============
\begin{frame}{Biomarqueurs Immunologiques Classiques}
  \textbf{Auto-anticorps standards [Katte 2023, Oram 2025]} :
  
  \begin{table}[h]
    \small
    \centering
    \begin{tabularx}{0.95\textwidth}{|l|c|c|X|}
      \toprule
      \textbf{Biomarqueur} & \textbf{Sensibilité} & \textbf{Spécificité} & \textbf{Performance Afrique} \\
      \midrule
      GADA (Anti-GAD65) & 60--70\% & Excellente & \textbf{Réduite} (formes idiopathiques) \\
      IA-2A (Tyrosine Phosphatase) & 50--60\% & Excellente & \textbf{Réduite} \\
      IAA (Anti-Insuline) & 40--50\% & Bonne & Faible \\
      ZnT8A (Zinc Transporter-8) & 50--60\% & Bonne & Complémentaire \\
      \bottomrule
    \end{tabularx}
  \end{table}
  
  \vspace{0.8em}
  
  \textbf{\textcolor{accent1}{Limites contexte africain}} :
  \begin{itemize}
    \item Coût prohibitif : \$50--150 USD/test
    \item Formes idiopathiques (30--40\%) : auto-anticorps négatifs
    \item Sensibilité diminuée populations africaines
    \item Inaccessibilité laboratoires spécialisés
  \end{itemize}
\end{frame}

% ============ SLIDE 9 ============
\begin{frame}{Biomarqueurs Métaboliques}
  \begin{columns}[T]
    \column{0.48\textwidth}
    \textbf{\color{accent2} Biomarqueurs métaboliques} :
    \begin{itemize}
      \item \textbf{Glycémie à jeun} : $\geq 7,0$ mmol/L
      \item \textbf{HGPO} : Hyperglycémie provoquée orale
      \item \textbf{HbA1c} : $\geq 6,5\%$ (glycation hémoglobine)
      \item \textbf{Peptide C} : Évaluation fonction $\beta$
      \item \textbf{HOMA-B} : Indice fonction cellules $\beta$
    \end{itemize}
    
    \column{0.48\textwidth}
    \textbf{\color{accent2} Stades TrialNet (3 stades)} :
    \begin{itemize}
      \item \textbf{Stade 1} : Autoimmunité sans dysglycémie
      \item \textbf{Stade 2} : Autoimmunité + dysglycémie asymptomatique
      \item \textbf{Stade 3} : Hyperglycémie symptomatique
    \end{itemize}
  \end{columns}
  
  \vspace{1em}
  
  \textbf{\textcolor{accent1}{Limites pour détection pré-symptomatique}} :
  \begin{itemize}
    \item Stades 1--2 difficiles à détecter sans biomarqueurs avancés
    \item Coût HGPO et HbA1c encore élevé en contexte africain
    \item Détection souvent au stade 3 (symptomatique)
  \end{itemize}
\end{frame}

% ============ SLIDE 10 ============
\begin{frame}{Biomarqueurs Génétiques et Scores de Risque}
  \begin{columns}[T]
    \column{0.48\textwidth}
    \textbf{\color{accent2} Gènes HLA classiques} :
    \begin{itemize}
      \item \textbf{HLA-DR/DQ} : 40--50\% risque total
      \item Haplotypes DQ2/DQ8 : majeurs Occident
      \item Haplotypes africains distincts
      \item Fréquences alléliques différentes
      \item Divergence évolutive millénaires
    \end{itemize}
    
    \column{0.48\textwidth}
    \textbf{\color{accent1} Gènes émergents (GWAS)} :
    \begin{itemize}
      \item \textbf{ANP32A-IT1} : ARN non-codant, régulation immune
      \item \textbf{ESCO2} : Apoptose, cycle cellulaire
      \item \textbf{NBPF1} : Expression pancréatique
      \item Identifiés par GWAS récents
      \item Rôle fonctionnel émergent
    \end{itemize}
  \end{columns}
  
  \vspace{1em}
  
  \textbf{\color{accent2} Genetic Risk Scores (GRS)} :
  \begin{itemize}
    \item AUC 0,85--0,92 populations occidentales
    \item \textbf{\textcolor{accent1}{Chute performance Afrique}} : AUC 0,75--0,80 [Oram 2025]
  \end{itemize}
\end{frame}

% ============ SLIDE 11 ============
\begin{frame}{Limites des Approches Diagnostiques en Contexte Africain}
  \begin{tcolorbox}[colback=accent1!15, title=\textbf{\small Barrières systémiques majeures}]
    \begin{itemize}
      \item \textbf{Coût tests immunologiques} : \$50--150 USD/test (auto-anticorps)
      \item \textbf{Coût tests génétiques} : \$50--200 USD/échantillon (NGS)
      \item \textbf{Manque laboratoires spécialisés} : Quasi-inexistants hors grandes villes
      \item \textbf{Personnel formé} : Diabétologie pédiatrique insuffisante
      \item \textbf{Infrastructure} : Électricité, chaîne du froid, réactifs
    \end{itemize}
  \end{tcolorbox}
  
  \vspace{0.8em}
  
  \textbf{\textcolor{accent1}{Hétérogénéité génétique africaine}} :
  \begin{itemize}
    \item Modèles occidentaux entraînés sur populations caucasiennes
    \item \textbf{Dataset shift} : $P_{\text{train}}(X, Y) \neq P_{\text{test}}(X, Y)$
    \item Perte de performance prédictive majeure
    \item Formes idiopathiques sous-représentées datasets occidentaux
  \end{itemize}
  
  \vspace{0.6em}
  
  \textbf{\textcolor{accent2}{Conclusion}} : Nécessité approche plus prédictive, accessible et adaptée
\end{frame}

% ============================================================================
% SECTION 3: APPRENTISSAGE AUTOMATIQUE POUR LA DÉTECTION DU DT1
% ============================================================================
\section{Apprentissage automatique pour la détection du DT1 : état de l'art}

% ============ SLIDE 12 ============
\begin{frame}{Problème de Classification Binaire en Médecine}
  \textbf{Formulation du problème} :
  
  \vspace{0.6em}
  
  \begin{equation*}
    \text{Classification binaire : } f(X) \rightarrow Y \in \{\text{Sain}, \text{À risque DT1}\}
  \end{equation*}
  
  \vspace{0.6em}
  
  \begin{columns}[T]
    \column{0.48\textwidth}
    \textbf{\color{accent2} Types de variables (X)} :
    \begin{itemize}
      \item \textbf{Cliniques} : âge, sexe, IMC, antécédents familiaux
      \item \textbf{Métaboliques} : glycémie, HbA1c, peptide C
      \item \textbf{Génétiques} : ANP32A-IT1, ESCO2, NBPF1, HLA
      \item Variables mixtes : 20--50 features typiques
    \end{itemize}
    
    \column{0.48\textwidth}
    \textbf{\color{accent2} Séparation données} :
    \begin{itemize}
      \item \textbf{Train} : 60--70\% (entraînement modèle)
      \item \textbf{Validation} : 15--20\% (optimisation hyperparamètres)
      \item \textbf{Test} : 15--20\% (évaluation finale)
      \item \textbf{Stratification} : Préserver ratio DT1+/DT1-
    \end{itemize}
  \end{columns}
\end{frame}

% ============ SLIDE 13 ============
\begin{frame}{Algorithmes Adaptés aux Données Biomédicales Tabulaires}
  \begin{table}[h]
    \tiny
    \centering
    \begin{tabularx}{\textwidth}{|l|c|c|c|X|}
      \toprule
      \textbf{Algorithme} & \textbf{Performance} & \textbf{Interprétabilité} & \textbf{Scalabilité} & \textbf{Caractéristiques} \\
      \midrule
      Régression Logistique & Moyenne & \textbf{Excellente} & Excellente & Baseline interprétable, coefficients linéaires \\
      k-NN & Bonne & Moyenne & Faible & Basé proximité, simple mais sensible bruit \\
      Random Forest & \textbf{Très bonne} & Bonne & Bonne & Ensemble arbres, feature importance, robuste \\
      SVM (RBF) & Bonne & Faible & Moyenne & Kernel trick, efficace haute dimension \\
      XGBoost & \textbf{Excellente} & Moyenne & Moyenne & Gradient boosting, état de l'art tabulaire \\
      LightGBM & \textbf{Excellente} & Moyenne & Bonne & Gradient boosting, plus rapide XGBoost \\
      \bottomrule
    \end{tabularx}
  \end{table}
  
  \vspace{0.6em}
  
  \textbf{\textcolor{accent2}{Choix projet : Random Forest + XGBoost}} :
  \begin{itemize}
    \item Performances supérieures sur données tabulaires
    \item Feature importance intégrée (interprétabilité)
    \item Robustesse au bruit et overfitting
  \end{itemize}
\end{frame}

% ============ SLIDE 14 ============
\begin{frame}{Déséquilibre de Classes et Métriques Cliniques}
  \textbf{\textcolor{accent1}{Problème critique}} : Déséquilibre typique 85--95\% sains, 5--15\% DT1
  
  \vspace{0.6em}
  
  \begin{columns}[T]
    \column{0.48\textwidth}
    \textbf{\color{accent1} Illusion Accuracy} :
    
    Modèle naïf « toujours sain » :
    
    \begin{itemize}
      \item \textbf{Accuracy = 90\%} ✓
      \item \textbf{Recall = 0\%} ✗ \textit{(catastrophe!)}
      \item Tous malades manqués
      \item Modèle inutile cliniquement
    \end{itemize}
    
    \column{0.48\textwidth}
    \textbf{\color{accent2} Métriques adaptées} :
    \begin{itemize}
      \item \textbf{Recall (Sensibilité)} : $\frac{VP}{VP + FN}$ \textbf{PRIORITAIRE}
      \item \textbf{F1-Score} : $2 \times \frac{\text{Prec} \times \text{Rec}}{\text{Prec} + \text{Rec}}$
      \item \textbf{AUC-ROC} : Indépendant seuil
      \item \textbf{AUC-PR} : Précision-Recall, adapté déséquilibre
    \end{itemize}
  \end{columns}
  
  \vspace{1em}
  
  \textbf{Justification clinique} : Manquer malade (FN) >> Diagnostiquer à tort sain (FP)
\end{frame}

% ============ SLIDE 15 ============
\begin{frame}{Rôle du Recall et F1-Score en Contexte Médical}
  \begin{columns}[T]
    \column{0.48\textwidth}
    \textbf{\color{accent2} Coûts asymétriques} :
    \begin{itemize}
      \item \textbf{Faux Négatif (FN)} : Patient DT1 non détecté
      \begin{itemize}
        \item Retard diagnostic
        \item Acidocétose potentiellement fatale
        \item Complications irréversibles
        \item \textbf{\textcolor{accent1}{Coût : tragédie}}
      \end{itemize}
      \vspace{0.4em}
      \item \textbf{Faux Positif (FP)} : Patient sain diagnostiqué DT1
      \begin{itemize}
        \item Investigation complémentaire
        \item Tests de confirmation
        \item Anxiété temporaire
        \item \textbf{Coût : acceptable}
      \end{itemize}
    \end{itemize}
    
    \column{0.48\textwidth}
    \textbf{\color{accent2} Solutions techniques} :
    \begin{itemize}
      \item \textbf{SMOTE} : Synthetic Minority Over-sampling Technique
      \item \textbf{Class weights} : Pénaliser erreurs classe minoritaire
      \item \textbf{Sous-échantillonnage} : Réduire classe majoritaire
      \item \textbf{Seuil de décision} : Ajuster pour maximiser Recall
    \end{itemize}
  \end{columns}
  
  \vspace{1em}
  
  \textbf{\textcolor{accent1}{Objectif clinique}} : Recall $\geq 90\%$ (détecter 9/10 malades minimum)
\end{frame}

% ============ SLIDE 16 ============
\begin{frame}{Interprétabilité : SHAP et LIME}
  \begin{columns}[T]
    \column{0.48\textwidth}
    \textbf{\color{accent2} SHAP (Lundberg 2017)} :
    \begin{itemize}
      \item Théorie des jeux (Shapley values)
      \item Contribution équitable de chaque feature
      \item \textbf{État de l'art 2020--2025}
      \item Explications globale + locale
      \item Robustesse théorique garantie
      \item Visualisations : waterfall, summary plot
    \end{itemize}
    
    \column{0.48\textwidth}
    \textbf{\color{accent2} LIME (Ribeiro 2016)} :
    \begin{itemize}
      \item Local Interpretable Model-agnostic Explanations
      \item Approximation locale par modèle simple
      \item Plus rapide que SHAP
      \item Moins robuste théoriquement
      \item Bon compromis vitesse/clarté
      \item Agnostique au modèle
    \end{itemize}
  \end{columns}
  
  \vspace{1em}
  
  \textbf{\textcolor{accent1}{Importance clinique}} :
  \begin{itemize}
    \item Confiance cliniciens nécessite transparence
    \item Régulation : RGPD (art. 22), EU AI Act 2026 $\Rightarrow$ explicabilité obligatoire
    \item Identifier biomarqueurs les plus discriminants
  \end{itemize}
\end{frame}

% ============ SLIDE 17 ============
\begin{frame}{Dataset Shift et Limites des Travaux Existants}
  \textbf{Problème critique : Dataset shift Occident $\rightarrow$ Afrique} :
  
  \vspace{0.6em}
  
  \begin{equation*}
    \boxed{P_{\text{train}}(X, Y) \neq P_{\text{test}}(X, Y)}
  \end{equation*}
  
  \vspace{0.6em}
  
  \begin{columns}[T]
    \column{0.48\textwidth}
    \textbf{\small Entraînement (Occident)} :
    \begin{itemize}
      \item >90\% cohortes caucasiennes
      \item Formes idiopathiques 5--10\%
      \item Haplotypes HLA DQ2/DQ8 majeurs
      \item Infrastructure diagnostique complète
    \end{itemize}
    
    \column{0.48\textwidth}
    \textbf{\small Déploiement (Afrique)} :
    \begin{itemize}
      \item Génétique africaine distincte
      \item Formes idiopathiques 30--40\%
      \item Haplotypes HLA différents
      \item Ressources diagnostiques limitées
    \end{itemize}
  \end{columns}
  
  \vspace{0.8em}
  
  \textbf{\textcolor{accent1}{Impact empirique}} : AUC 0,90--0,92 (Occident) $\rightarrow$ 0,75--0,80 (Afrique) [Oram 2025]
  
  \vspace{0.6em}
  
  \textbf{\textcolor{accent1}{Rareté critique}} : \textbf{ZÉRO étude ML} développée spécifiquement sur données camerounaises/africaines DT1
\end{frame}

% ============================================================================
% SECTION 4: MÉTHODOLOGIE COMPUTATIONNELLE DU PROJET
% ============================================================================
\section{Méthodologie computationnelle du projet}

% ============ SLIDE 18 ============
\begin{frame}{Données et Prétraitement}
  \begin{columns}[T]
    \column{0.48\textwidth}
    \textbf{\color{accent2} Sources de données} :
    \begin{itemize}
      \item \textbf{Cohortes hospitalières camerounaises}
      \item Hôpitaux de référence (Yaoundé, Douala)
      \item Données cliniques + métaboliques
      \item Biomarqueurs génétiques (ANP32A-IT1, ESCO2, NBPF1, HLA)
      \item \textit{Note : Données synthétiques si réelles indisponibles}
    \end{itemize}
    
    \column{0.48\textwidth}
    \textbf{\color{accent2} Critères inclusion/exclusion} :
    \begin{itemize}
      \item \textbf{Inclusion} : Âge $<40$ ans, diagnostic confirmé
      \item \textbf{Exclusion} : Données manquantes $>30\%$, DT2 confirmé
      \item Consentement éclairé (éthique)
      \item Taille cible : $n \geq 500$ patients
    \end{itemize}
  \end{columns}
  
  \vspace{1em}
  
  \textbf{\color{accent2} Pipeline de prétraitement} :
  \begin{enumerate}
    \item \textbf{Nettoyage} : Suppression doublons, correction erreurs saisie
    \item \textbf{Valeurs manquantes} : Imputation médiane (si $<20\%$) ou suppression
    \item \textbf{Encodage} : Variables catégorielles (One-Hot Encoding)
    \item \textbf{Normalisation} : StandardScaler (moyenne=0, écart-type=1)
  \end{enumerate}
\end{frame}

% ============ SLIDE 19 ============
\begin{frame}{Pipeline d'Apprentissage : Découpage et Déséquilibre}
  \textbf{\color{accent2} Découpage des données} :
  
  \vspace{0.6em}
  
  \begin{columns}[T]
    \column{0.48\textwidth}
    \textbf{Stratégie retenue} :
    \begin{itemize}
      \item \textbf{Train} : 60\% (entraînement)
      \item \textbf{Validation} : 20\% (optimisation hyperparamètres)
      \item \textbf{Test} : 20\% (évaluation finale)
      \item \textbf{Stratification} : Préserver ratio DT1+/DT1-
      \item \textbf{Validation croisée} : k-fold (k=5) pour robustesse
    \end{itemize}
    
    \column{0.48\textwidth}
    \textbf{Gestion du déséquilibre} :
    \begin{itemize}
      \item \textbf{SMOTE} : Génération synthétique instances minoritaires
      \item \textbf{Class weights} : Pénaliser erreurs DT1+ (classe minoritaire)
      \item \textbf{Stratification} : Obligatoire tous splits
      \item \textbf{Seuil ajusté} : Maximiser Recall
    \end{itemize}
  \end{columns}
  
  \vspace{1em}
  
  \textbf{\textcolor{accent1}{Choix SMOTE}} : Évite perte information (vs sous-échantillonnage) tout en équilibrant classes
\end{frame}

% ============ SLIDE 20 ============
\begin{frame}{Modèles Entraînés et Optimisation}
  \textbf{\color{accent2} Modèles effectivement entraînés} :
  
  \vspace{0.6em}
  
  \begin{columns}[T]
    \column{0.48\textwidth}
    \textbf{Random Forest} :
    \begin{itemize}
      \item \textbf{Hyperparamètres optimisés} :
      \begin{itemize}
        \item n\_estimators : \{100, 200, 300\}
        \item max\_depth : \{10, 20, 30, None\}
        \item min\_samples\_split : \{2, 5, 10\}
      \end{itemize}
      \item Feature importance intégrée
      \item Robuste au bruit
    \end{itemize}
    
    \column{0.48\textwidth}
    \textbf{XGBoost} :
    \begin{itemize}
      \item \textbf{Hyperparamètres optimisés} :
      \begin{itemize}
        \item n\_estimators : \{100, 200, 300\}
        \item max\_depth : \{3, 5, 7\}
        \item learning\_rate : \{0.01, 0.1, 0.3\}
      \end{itemize}
      \item Gradient boosting optimisé
      \item État de l'art tabulaire
    \end{itemize}
  \end{columns}
  
  \vspace{1em}
  
  \textbf{\color{accent2} Stratégie d'optimisation} :
  \begin{itemize}
    \item \textbf{GridSearchCV} : Recherche exhaustive hyperparamètres
    \item \textbf{Validation croisée 5-fold} : Robustesse évaluation
    \item \textbf{Métrique optimisée} : F1-Score (équilibre Precision/Recall)
  \end{itemize}
\end{frame}

% ============ SLIDE 21 ============
\begin{frame}{Évaluation et Métriques Retenues}
  \textbf{\color{accent2} Métriques d'évaluation prioritaires} :
  
  \vspace{0.6em}
  
  \begin{table}[h]
    \small
    \centering
    \begin{tabularx}{0.90\textwidth}{|l|X|c|}
      \toprule
      \textbf{Métrique} & \textbf{Signification clinique} & \textbf{Priorité} \\
      \midrule
      \textbf{Recall} & \% vrais DT1 détectés (ne pas manquer malades) & \textbf{Maximale} \\
      \textbf{F1-Score} & Équilibre Precision/Recall & Élevée \\
      \textbf{AUC-ROC} & Capacité discrimination globale & Élevée \\
      \textbf{AUC-PR} & Précision-Recall (adapté déséquilibre) & Moyenne \\
      \textbf{Accuracy} & \% prédictions correctes & Faible \\
      \bottomrule
    \end{tabularx}
  \end{table}
  
  \vspace{0.8em}
  
  \textbf{\color{accent2} Critères de choix du modèle final} :
  \begin{enumerate}
    \item \textbf{Recall $\geq 90\%$} : Priorité absolue (détection malades)
    \item \textbf{AUC-ROC $\geq 0.85$} : Performance globale suffisante
    \item \textbf{F1-Score maximal} : Parmi modèles satisfaisant critères 1--2
    \item \textbf{Interprétabilité} : Départage ex aequo (Random Forest favori)
  \end{enumerate}
\end{frame}

% ============ SLIDE 22 ============
\begin{frame}{Interprétabilité : Analyse SHAP}
  \textbf{\color{accent2} Analyse SHAP (SHapley Additive exPlanations)} :
  
  \vspace{0.6em}
  
  \begin{columns}[T]
    \column{0.48\textwidth}
    \textbf{Importance globale} :
    \begin{itemize}
      \item \textbf{Summary plot} : Ranking biomarqueurs
      \item Identifier top 5--10 features discriminantes
      \item Comparaison ANP32A-IT1, ESCO2, NBPF1 vs HLA
      \item Validation cohérence biologique
    \end{itemize}
    
    \column{0.48\textwidth}
    \textbf{Importance locale (par patient)} :
    \begin{itemize}
      \item \textbf{Waterfall plot} : Explication individuelle
      \item Contribution chaque biomarqueur à prédiction
      \item Transparence pour clinicien
      \item Confiance décision
    \end{itemize}
  \end{columns}
  
  \vspace{1em}
  
  \textbf{\color{accent2} Visualisations pour cliniciens} :
  \begin{itemize}
    \item \textbf{Top features} : Tableau biomarqueurs ordonnés par importance
    \item \textbf{Explication patient} : Rapport individuel (PDF) avec SHAP waterfall
    \item \textbf{Interface simple} : Application Streamlit/Flask pour déploiement
  \end{itemize}
\end{frame}

% ============ SLIDE 23 ============
\begin{frame}{Synthèse de la Méthodologie}
  \begin{tcolorbox}[colback=accent2!10, title=\textbf{\small Pipeline computationnel complet}]
    \textbf{1. Données} : Cohortes camerounaises, biomarqueurs génétiques (ANP32A-IT1, ESCO2, NBPF1, HLA) + cliniques
    
    \textbf{2. Prétraitement} : Nettoyage, imputation, normalisation, stratification
    
    \textbf{3. Gestion déséquilibre} : SMOTE + class weights
    
    \textbf{4. Modèles} : Random Forest + XGBoost (optimisation GridSearchCV)
    
    \textbf{5. Évaluation} : Recall $\geq 90\%$ (priorité), AUC-ROC $\geq 0.85$, F1-Score
    
    \textbf{6. Interprétabilité} : SHAP (globale + locale), visualisations cliniciens
  \end{tcolorbox}
  
  \vspace{0.8em}
  
  \textbf{\textcolor{accent1}{Originalité scientifique}} :
  \begin{itemize}
    \item \textbf{PREMIÈRE} approche ML DT1 données africaines
    \item Adaptation explicite contexte camerounais (formes idiopathiques, dataset shift)
    \item Rigueur méthodologique : validation croisée, métriques adaptées
  \end{itemize}
\end{frame}

% ============ SLIDE 24 ============
\begin{frame}{Vision et Impact Attendu}
  \begin{tcolorbox}[colback=accent2!20, fonttitle=\bfseries, title=Contribution Scientifique Majeure]
    Ce projet établit les \textbf{fondements} pour systèmes de dépistage DT1 accessibles, coût-efficaces et culturellement adaptés, réduisant la mortalité évitable en Afrique subsaharienne.
  \end{tcolorbox}
  
  \vspace{0.8em}
  
  \textbf{\color{accent2} Impact scientifique} :
  \begin{itemize}
    \item Publication originale : \textbf{PREMIÈRE étude ML DT1 Afrique}
    \item Protocole reproductible et open-source
    \item Données et code partagés (sous éthique appropriée)
  \end{itemize}
  
  \vspace{0.6em}
  
  \textbf{\color{accent2} Impact clinique et santé publique} :
  \begin{itemize}
    \item Réduction mortalité 1ère année (objectif : 25--30\% $\rightarrow$ $<15\%$)
    \item Détection pré-symptomatique pour prévention acidocétose
    \item Outil accessible contexte ressources limitées
  \end{itemize}
  
  \vspace{0.6em}
  
  \textbf{\textcolor{accent1}{Enjeu éthique}} : Éviter exacerbation inégalités sanitaires par modèles inadaptés
\end{frame}

% ============ SLIDE 25: MERCI ============
\begin{frame}[plain]
  \Huge
  \centering
  \textbf{Merci de votre attention !}
  
  \vspace{2em}
  
  \Large
  \textit{Détection précoce du Diabète de Type 1 au Cameroun \\ par Apprentissage Automatique}
  
  \vspace{1em}
  
  \normalsize
  \textit{Utilisation de biomarqueurs génétiques émergents \\ (ANP32A-IT1, ESCO2, NBPF1)}
  
  \vspace{2em}
  
  \normalsize
  Supervision : Dr. TCHAPET NJAFA Jean-Pierre \\
  Conseiller scientifique : Prof. NANA Engo
\end{frame}


\end{document}
